\section{\texorpdfstring{$\overline\partial_{(j,J)}$}{dbar(j,J)} 的线性化}
\label{sec:linearizations-dbar-jj}
\setcounter{equation}{4}

设 $\Sigma$ 是亏格为 $g$ 的紧曲面，并像以前一样，以 $\mathcal J_\Sigma$ 表示 $\Sigma$ 上复结构 $j$ 的集合。

首先固定 $\Sigma$ 上的一个复结构 $j$。

\begin{proposition}\label{prop:dbar-fredholm-section}
设 $J$ 是 $(M,\omega)$ 上任意相容殆复结构。如果在
$T_u\mathcal F$ 上赋予 $W^{k,p}$ 拓扑，在
$\mathcal H^{(0,1)}_{(j,J);u}$ 上赋予 $W^{k-1,p}$ 拓扑，其中
$k-2/p>0$，那么赋值
\[
 \overline\partial_{(j,J)}:\mathcal F\longrightarrow
 \mathcal H^{(0,1)}_{(j,J)}
\]
定义一个弗雷德霍姆截面。
\end{proposition}

换言之，线性化映射
\[
 D_u\overline\partial_{(j,J)}:
 T_u\mathcal F\longrightarrow\mathcal H^{(0,1)}_{(j,J);u}
\]
是弗雷德霍姆算子。本节余下部分将证明这个命题，并计算该弗雷德霍姆算子的指标。

回忆
\[
 T_u\mathcal F=\Omega^0(u^*TM),
 \qquad
 \mathcal H^{(0,1)}_{(j,J);u}
 =\Omega^{(0,1)}_{(j,J)}(u^*TM).
\]
为此要显式计算 $D_u\overline\partial_{(j,J)}$。计算中将使用殆凯勒结构
$(M,\omega,J)$ 的一个殆埃尔米特联络（例如典范联络），也就是满足
\[
 \nabla\omega=0=\nabla J
\]
的联络。

这个联络的挠率未必为零；一般而言，除非 $J$ 可积，否则它不同于 $g_J$ 的列维--奇维塔联络。

假设 $d(x_0,x_1)<\operatorname{inj}(M,\nabla)$，从而关于联络 $\nabla$ 存在一条从 $x_0$ 到 $x_1$ 的唯一短测地线。以
\[
 \Pi_{x_0}^{x_1}:T_{x_0}M\longrightarrow T_{x_1}M
\]
表示沿该唯一短测地线的平行移动。

我们先求截面 $d$ 的（协变）线性化的显式表达式。$TM$ 上的联络 $\nabla$ 也通过沿短测地线的平行移动在向量丛
$\mathcal H\to\mathcal F$ 上给出一个联络 $\nabla$，具体如下。

设 $X$ 是 $M$ 上的向量场，$v\in T_{x_0}M$。令
$\gamma:(-\epsilon,\epsilon)\to M$ 是满足
$\gamma(0)=x_0$、$\gamma'(0)=v$ 的唯一测地线。则有公式
\[
 \nabla_vX=
 \left.\frac d{ds}\right|_{s=0}
 \left(\Pi_{x_0}^{\gamma(s)}\right)^{-1}
 \bigl(X(\gamma(s))\bigr)\in T_{x_0}M.
\]

对映射 $u:\Sigma\to M$，以
$\nabla_{du}:=u^*\nabla$ 表示向量丛 $u^*TM\to\Sigma$ 上 $\nabla$ 的拉回联络。

给定任意截面 $\Upsilon:\mathcal F\to\mathcal H$（例如 $du$）和
$\xi\in T_u\mathcal F$，定义协变导数对 $\xi$ 在 $z$ 处的取值为
\[
 \nabla_{du}\Upsilon(\xi)(z)
 =\left.\frac d{ds}\right|_{s=0}
 \left(\Pi_{u(z)}^{u_\xi^s(z)}\right)^{-1}
 \bigl(\Upsilon(u_\xi^s(z))\bigr),
\]
其中 $z\in\Sigma$，并且
$u_\xi^s(z):=\exp_{u(z)}s(\xi(z))$。下面显式计算协变线性化
$D_ud:T_u\mathcal F\to\mathcal H_u$。回忆 $du$ 是 $\Sigma$ 上取值于
$u^*TM$ 的 $1$-形式，因此尝试对 $a\in T_z\Sigma$ 计算
$D_ud(\xi)(a)$。

首先，取曲线 $\lambda:(-\epsilon,\epsilon)\to\Sigma$，满足
$\lambda(0)=z$、$\lambda'(0)=a$，并考虑单参数曲线族
\[
 \Gamma:(-\epsilon,\epsilon)\times(-\epsilon,\epsilon)
 \longrightarrow M,
 \qquad
 (t,s)\longmapsto\exp_{u(\lambda(t))}(s\xi).
\]
注意
$\Gamma(0,s)=\exp_zs\xi=u_\xi^s(z)$，
$\Gamma(t,0)=u(\lambda(t))$。因此
\[
 \left.\frac{\partial\Gamma}{\partial t}\right|_{t=0,s=0}=du(a),
 \qquad
 \left.\frac{\partial\Gamma}{\partial s}\right|_{s=0,t=0}=\xi(u(z)).
\]
回忆下述挠率公式：
\[
 \left.\nabla_s\frac{\partial\Gamma}{\partial s}
 \right|_{(t,s)=(0,0)}
 -\left.\nabla_t\frac{\partial\Gamma}{\partial s}
 \right|_{(t,s)=(0,0)}
 =T\left(
 \left.\frac{\partial\Gamma}{\partial s}\right|_{(s,t)=(0,0)},
 \left.\frac{\partial\Gamma}{\partial t}\right|_{(s,t)=(0,0)}
 \right)
 =T(\xi,du(a)).
\]

现在计算 $D_u(d)(\xi)(a)$：
\begin{align*}
 D_u(d)(\xi)(a)
 &=\left.\frac D{\partial s}\right|_{(s,t)=(0,0)}
 d(u^s)\left(\frac{d\gamma}{dt}\right)
 =\left.\frac D{\partial s}\frac{\partial\Gamma}{\partial t}
 \right|_{(s,t)=(0,0)}\\
 &=\left.\frac D{\partial t}\frac{\partial\Gamma}{\partial s}
 \right|_{(s,t)=(0,0)}+T(\xi,du(a))\\
 &=\left.\frac D{\partial t}\right|_{t=0}\xi+T(\xi,du(a))
 =\nabla_{du(a)}\xi+T(\xi,du(a)).
\end{align*}

现在以 $\widetilde T_{du}$ 表示 $\Sigma$ 上取值于
$\operatorname{End}(u^*TM)$ 的 $1$-形式，它定义为
\[
 \widetilde T_{du}(\xi)(a)=T(\xi,du(a)),
 \qquad a\in T\Sigma,\ \xi\in u^*TM.
\]
用这些记号，得到 $d$ 的（协变）线性化公式
\begin{equation}
 D_ud=\nabla_{du}+\widetilde T_{du}.
 \label{eq:covariant-linearization-d}
\end{equation}
于是由联络 $\nabla$ 的殆凯勒性质，得到
\[
 D_u\overline\partial_{(j,J)}
 =(\nabla_{du}+\widetilde T_{du})^{(0,1)}_{(j,J)}.
\]
更显式地，可写成
\begin{equation}
 \begin{split}
 D_u\overline\partial_{(j,J)}(\xi)(a)
 ={}&\frac12\bigl(\nabla_{du(a)}\xi
 +J\cdot\nabla_{du(j,a)}\xi\bigr)\\
 &+\frac12\bigl(T(\xi,du(a))
 +J\cdot T(\xi,du(j,a))\bigr).
 \end{split}
 \label{eq:explicit-linearization-dbar-jj}
\end{equation}
将以上计算概括为下面的命题。

\begin{proposition}\label{prop:linearization-dbar-elliptic}
设 $\nabla$ 是 $(M,\omega,J)$ 的一个殆埃尔米特联络。固定
$\Sigma$ 上的 $j$，并把 $\overline\partial_{(j,J)}$ 看成丛
$\mathcal H^{(0,1)}_{(j,J)}\to\mathcal F$ 的截面。则截面
$\overline\partial_{(j,J)}$ 在光滑映射 $u:\Sigma\to M$ 处的协变线性化
\[
 D_u\overline\partial_{(j,J)}:
 \Omega^0(u^*TM)\longrightarrow
 \Omega^{(0,1)}_{(j,J)}(u^*TM)
\]
是一个椭圆的一阶微分算子。
\end{proposition}

\begin{proof}
由\eqref{eq:explicit-linearization-dbar-jj}，
$D_u\overline\partial_{(j,J)}$ 是一阶微分算子，其符号由函数
\[
 \sigma\bigl(D_u\overline\partial_{(j,J)}\bigr):
 T^*\Sigma\setminus o_{T^*\Sigma}
 \longrightarrow
 \operatorname{Hom}_{\mathbb R}
 \bigl(\pi^*(u^*TM),\pi^*(\Lambda^1(u^*TM))\bigr)
\]
给出，具体为
\[
 \alpha_z\longmapsto
 \frac{\alpha_z\otimes\operatorname{Id}
 +\alpha_z\circ j(z)\otimes J(u(z))}{2}.
\]
从 $\alpha_z$ 的像的这一表达式容易验证：对任意 $\alpha_z\ne0$，它都是同构。这就证明该算子是椭圆的。
\end{proof}

需要指出，尽管 $D_u\overline\partial_{(j,J)}$ 的表达式看起来依赖于联络的选择，当 $u$ 是 $(j,J)$-全纯的，即
$\overline\partial_{(j,J)}u=0$ 时，它实际上与联络的选择无关。

结合这一点与黎曼--罗赫公式（或阿蒂亚--辛格指标公式），立即得到下面的定理。

\begin{theorem}\label{thm:index-linearized-dbar}
假设 $(\Sigma,j)$ 是亏格为 $g$ 的闭黎曼曲面，且
$[u]=\beta\in H_2(M,\mathbb Z)$。则
\[
 D_u\overline\partial_{(j,J)}:
 T_u\mathcal F^{(k,p)}\longrightarrow\mathcal H_{(j,J)}^{(k-1,p)}
\]
是弗雷德霍姆算子，其指标为
\[
 \operatorname{Index}D_u\overline\partial_{(j,J)}
 =2c_1(\beta)+n\chi(\Sigma)
 =2c_1(\beta)+2n(1-g).
\]
因此，$\overline\partial_{(j,J)}$ 是巴拿赫丛
$\mathcal H_{(j,J)}^{(k-1,p)}\to\mathcal F^{(k,p)}$ 的弗雷德霍姆截面。
\end{theorem}

现在固定同调类 $\beta\in H_2(M,\mathbb Z)$。考虑 $\mathcal F^{k,p}$ 的子集
\[
 \mathcal F^{k,p}(\beta)=
 \{u\in\mathcal F^{k,p}\mid [u]=\beta\}
\]
以及零点集
\[
 \widetilde{\mathcal M}^{k,p}(J;\beta)=
 \{u\in\mathcal F^{k,p}(\beta)
 \mid\overline\partial_{(j,J)}u=0\},
\]
其中 $k\geq1$、$p>2$。由椭圆正则性，
$\widetilde{\mathcal M}^{k,p}(J;\beta)$ 的任意元素都是光滑的，因而该集合不依赖于 $k,p$。把这个公共空间记作
$\widetilde{\mathcal M}(J;\beta)$，并将它看成 $C^\infty(\Sigma,M)$ 的子集。再次由椭圆正则性与自举论证，
$\widetilde{\mathcal M}(J;\beta)$ 上由 $W^{k,p}$ 诱导的拓扑全都等价于从
$C^\infty(\Sigma,M)$ 诱导的光滑拓扑。
